\documentclass[11pt,a4paper]{ctexart}
\usepackage[left=1.8cm,right=1.8cm,top=1.65cm,bottom=1.65cm]{geometry}
\usepackage{enumitem}
\usepackage{xcolor}
\usepackage{hyperref}
\usepackage{fancyhdr}
\usepackage{tabularx}
\usepackage{array}
\usepackage{setspace}
\usepackage{amssymb}

\setmainfont{Times New Roman}
\setsansfont{Arial}
\hypersetup{colorlinks=true,linkcolor=black,urlcolor=blue}
\definecolor{accent}{HTML}{315D7A}
\definecolor{light}{HTML}{EDF3F6}
\definecolor{muted}{HTML}{555555}

\pagestyle{fancy}
\fancyhf{}
\lhead{\small 华东理工大学计算机方向线上交流准备}
\rhead{\small 刘鹏}
\cfoot{\small \thepage}
\setlength{\headheight}{14pt}
\setlength{\parindent}{0pt}
\setlength{\parskip}{0.35em}
\setlist[itemize]{leftmargin=1.7em,itemsep=0.18em,topsep=0.25em}
\setlist[enumerate]{leftmargin=1.9em,itemsep=0.3em,topsep=0.25em}
\renewcommand{\arraystretch}{1.25}

\newcommand{\sectionbar}[1]{%
  \vspace{0.35em}
  {\large\bfseries\sffamily\color{accent} #1}\par
  \vspace{-0.35em}\color{accent}\rule{\textwidth}{0.7pt}\color{black}\vspace{0.2em}
}
\newcommand{\prompt}[1]{\textbf{\color{accent}#1}}
\newcommand{\note}[1]{%
  \begin{center}
  \colorbox{light}{\parbox{0.94\textwidth}{#1}}
  \end{center}
}
\newcommand{\blankline}{\vspace{0.6em}\hrulefill\par}

\begin{document}

\begin{center}
  {\LARGE\bfseries\sffamily 华东理工大学计算机方向}\\[0.25em]
  {\LARGE\bfseries\sffamily 线上交流准备手册}\\[0.6em]
  {\large 刘鹏 \quad｜\quad 无 PPT 交流版}
\end{center}

\note{\textbf{交流目标：}让团队负责人在较短时间内确认三点：研究方向匹配、具备独立发展潜力、加入团队后能够形成明确增量。回答时先给结论，再给一至两个证据，单次回答尽量控制在 1--2 分钟。}

\sectionbar{一页速查}

\prompt{一句话定位}

围绕三维医学与生物图像分析，研究标签高效学习、实例分割和结构化推理，并形成“算法方法、数据基准、开源工具”相结合的研究体系。

\prompt{三条经历主线}
\begin{itemize}
  \item 博士阶段：半监督医学图像分类与多器官分割，代表工作发表于 JBHI、CBM 和 MICCAI。
  \item 博士后阶段：三维体电镜图像分析、细胞器实例分割、分割错误检测与人在环路校对。
  \item 工业经历：曾从事搜索与推荐算法研发，具备模型落地和工程协作经验。
\end{itemize}

\prompt{优先讲的两组成果}
\begin{itemize}
  \item \textbf{标签高效医学图像分析：}不平衡半监督分类、CVCL、CR-GLoCo。
  \item \textbf{体电镜与生物图像：}MitoEM 2.0、CellMap Challenge 第一名、CellAble，以及细胞核和细胞器分割合作。
\end{itemize}

\prompt{未来三条方向}
\begin{itemize}
  \item 三维医学与生物图像基础模型。
  \item 半监督、自监督及跨模态的标签高效学习。
  \item 融合形态与拓扑先验的可信分析和人在环路校对。
\end{itemize}

\prompt{希望传递的个人特点}
\begin{itemize}
  \item 研究主线连续：从医学影像扩展到体电镜，本质均是三维视觉、少标注学习和结构化分割。
  \item 能独立推进：可从问题定义、算法、实验评测推进到论文、代码和数据发布。
  \item 适合交叉合作：能够连接计算机视觉、智能医学和生命科学应用。
\end{itemize}

\prompt{交流结束前至少问}
\begin{itemize}
  \item 团队当前最希望新教师补充哪类研究能力？
  \item 新教师前三年的考核、招生和启动支持如何安排？
  \item 团队已有的医学、生物或其他交叉合作基础有哪些？
\end{itemize}

\newpage
\sectionbar{自我介绍}

\prompt{30 秒版本}
王老师您好，我叫刘鹏，硕士毕业于中科院，博士毕业于交大生物医学工程专业。目前是在美国波士顿学院计算机系做博士后。我的研究主线是三维医学与生物图像分析，主要结果标注成本高、结构关系复杂以及模型泛化能力不足的问题。

博士期间主要是研究半监督学习，目的是减少对于标注数据的需求，因为在医学图像领域，标注大量数据比较昂贵。博士后阶段主要研究的是三维体电镜图像。这里的数据规模大、实例密集、形态和拓扑学复杂，对算法的泛化能力和结构理解提出了更好的要求。

未来我希望围绕三维医学与生物图像的基础模型、标签高效与跨模态学习等方向继续研究。

老师您好，我叫刘鹏，博士毕业于上海交通大学生物医学工程专业，目前在美国波士顿学院计算机系从事博士后研究。我的研究主线是三维医学与生物图像分析，重点关注标签高效学习、细胞器实例分割和结构化推理。博士阶段主要研究半监督医学图像分类与多器官分割；博士后阶段进一步拓展到体电镜图像、数据基准和人在环路校对。未来希望在计算机学院形成算法、数据和工具协同发展的研究方向。

\prompt{约 3 分钟版本}

老师您好，我叫刘鹏。我的教育背景横跨地理信息、信号处理和生物医学工程，目前在美国波士顿学院计算机系开展博士后研究。虽然应用对象经历了遥感、医学影像和体电镜图像的变化，但我的研究主线比较连续，核心都是面向复杂二维和三维视觉数据，解决标注成本高、结构关系复杂以及模型泛化能力不足的问题。

博士期间，我主要研究标签高效的医学图像分析。在半监督分类方面，我针对类别不平衡问题提出了不确定性引导的虚拟对抗训练和批次核范数优化；在三维分割方面，我提出 CVCL 和 CR-GLoCo，分别从体素级对比学习、全局与局部上下文一致性的角度提升少标注条件下的多器官分割性能。相关工作发表于 JBHI、Computers in Biology and Medicine、MICCAI 和 Computerized Medical Imaging and Graphics。

博士后阶段，我把研究进一步延伸到三维体电镜图像。这里的数据规模大、实例密集、形态和拓扑结构复杂，对算法的泛化能力和结构理解提出了更高要求。我参与 MitoEM 2.0 数据基准建设，负责基线模型和评测流程；同时开展多细胞器分割、错误检测和自动校对研究，并获得 CellMap Segmentation Challenge 第一名。我也开发了 CellAble 标注与校对工具，希望把算法研究和实际数据处理流程连接起来。

未来我希望围绕三维医学与生物图像基础模型、标签高效与跨模态学习、结构先验与可信分析三个方向继续发展，形成“算法方法、数据基准、开源工具”相结合的研究体系。我认为这些工作既有计算机视觉和机器学习方面的方法问题，也适合与医学、生命科学等方向开展交叉合作。

\note{\textbf{表达提醒：}不用逐句背诵。记住“当前身份—博士工作—博士后工作—统一主线—未来计划”五个节点即可。对方若中途追问，应停止背稿，直接进入讨论。}

\sectionbar{两个代表工作故事}

\prompt{故事一：标签高效医学图像分析}

\textbf{问题：}医学图像标注依赖专家，而且常伴随类别不平衡、部分器官未标注和三维计算开销高等问题。

\textbf{做法：}围绕一致性学习与对比学习，分别从不确定性筛选、体素级表征学习及跨分辨率全局--局部协同三个角度提高未标注数据利用效率。

\textbf{证据：}CVCL 在五个公开 CT 数据集上取得五器官平均 DSC 92.5\%；CR-GLoCo 在 BCV、MMWHS 和 LA 三项任务分别达到 74.0\%、82.2\% 和 90.0\% DSC。

\textbf{个人贡献：}负责问题定义、方法设计、实验评测和论文撰写，并持续将博士阶段的工作发展为可扩展的研究方向。

\textbf{下一步：}从单任务半监督学习扩展到多任务基础模型、跨模态预训练和跨中心泛化。

\newpage
\prompt{故事二：体电镜与生物图像分析}

\textbf{问题：}体电镜数据体量大、实例密集且形态复杂，单纯追求像素精度不足以满足三维重建和生物定量分析需求。

\textbf{做法：}从数据基准、分割模型、错误检测和人在环路校对四个环节推进研究。参与 MitoEM 2.0 建设，负责基线模型和评测流程；开展多细胞器分割与结构错误校正；开发 CellAble 工具支持人工标注与模型结果校对。

\textbf{证据：}MitoEM 2.0 覆盖 8 个跨物种、组织及成像模态数据集，包含 28 个体数据和 8,000 余个专家校验线粒体实例；获得 CellMap Segmentation Challenge 第一名。

\textbf{个人贡献：}能够将算法研发延伸到数据组织、基准评测、工具实现和跨团队协作，而不仅是完成单一模型实验。

\textbf{下一步：}构建跨细胞器、跨成像模态的三维基础模型，研究融合形态和拓扑先验的实例分割及自动校对。

\sectionbar{未来三年计划}

\prompt{总体目标}

建立面向三维医学与生物图像的智能分析研究方向，形成稳定的论文产出、数据与软件资源，并逐步建设一支由研究生组成的小型团队。

\begin{tabularx}{\textwidth}{>{\bfseries}p{1.8cm} X}
第 1 年 & 完成现有博士后工作的整理与发表；搭建三维图像训练和评测平台；围绕标签高效三维分割形成一个可快速启动的课题；积极申请国家自然科学基金青年项目及校内启动项目。\\
第 2 年 & 推进跨模态或多任务预训练研究；与校内医学、生命科学或相关应用团队建立稳定合作；发布可复用的数据、代码或工具；培养研究生形成相对独立的子课题。\\
第 3 年 & 将基础模型、结构化推理和人在环路校对整合为较完整的研究体系；冲击高水平期刊和会议成果；为后续面上项目或交叉重点课题积累基础。\\
\end{tabularx}

\prompt{可能申请的项目}
\begin{itemize}
  \item 国家自然科学基金青年科学基金：标签高效三维医学/生物图像分析。
  \item 面向三维体电镜的结构先验实例分割与错误校正。
  \item 校内交叉项目：结合生命健康、生物医学或工业视觉场景。
\end{itemize}

\prompt{可以承担的课程}

程序设计、数据结构、机器学习、深度学习、计算机视觉、医学图像分析。交流时应说明可根据学院教学安排承担基础课程，不只强调专业选修课。

\newpage
\sectionbar{高频问题与回答框架}

\begingroup
\small

\prompt{1. 为什么考虑华东理工大学？}

建议从三个层面回答：一是计算机学科平台能够承载机器学习与计算机视觉研究；二是学校具有鲜明的理工和交叉应用背景，有利于发展智能医学与生物图像分析；三是上海具备医院、科研院所和产业合作资源。最后落到个人贡献：希望补充三维医学与生物图像、标签高效学习和体电镜分析方向。

\prompt{2. 从医学图像转到体电镜，研究是否过于分散？}

不是应用方向的简单切换。两类数据共享三维视觉、昂贵标注、复杂结构和跨域泛化等核心问题。医学图像阶段积累了半监督学习和多器官分割方法；体电镜阶段进一步研究实例级理解、拓扑结构和自动校对。统一主线是“面向少标注复杂三维图像的结构化学习”。

\prompt{3. 哪些工作是你独立主导的？}

选择两项回答，不要罗列全部论文。说明自己在问题定义、方法设计、实验实现、论文撰写和合作组织中的具体职责。对合作论文严格区分“负责”“共同完成”和“参与”，避免把团队成果全部表述为个人主导。

\prompt{4. 入职后第一年能做出什么？}

第一，整理并完成当前体电镜方向的在研工作；第二，复用现有代码和数据基础，快速启动标签高效三维分割课题；第三，申请青年基金和校内启动项目；第四，尽快参与研究生培养和课程教学。

\prompt{5. 你的研究如何与团队结合？}

先询问对方现有方向，再选择接口：若团队偏机器学习，可强调少样本、自监督和基础模型；若偏视觉，可强调三维分割、实例理解和结构化推理；若偏应用，可强调医学、生物数据、开源工具及跨团队协作能力。

\prompt{6. 你能承担哪些教学任务？}

已有程序设计课程助教经历，负责作业布置、上机指导、批改和答疑。可承担程序设计、数据结构等基础课，也可以逐步开设机器学习、计算机视觉或医学图像分析方向课程。

\prompt{7. 如何看待论文、项目和工程之间的关系？}

论文解决方法创新，项目提供真实问题和持续数据，工程工具负责把算法嵌入工作流程。自己的目标不是只做应用系统，也不是只追求单项指标，而是用真实问题推动可复用的方法研究。

\prompt{8. 医学图像和 EM 图像处理任务有什么不同？}

建议按“共同点--差异--迁移”回答。共同点是两者都面临少标注、复杂三维视觉、数据异质和结构化分割问题。差异在于：医学图像通常面向器官、组织和病灶，服务临床诊断、治疗规划和预后分析；EM 图像处于纳米尺度，更关注细胞和细胞器实例分割、三维重建、连接关系分析和自动校对。错误类型也不同，医学图像常见漏分、边界偏差和跨中心泛化问题，EM 中相邻实例粘连、断裂和拓扑错误更关键，一处错误可能影响整个三维实例。

收束时强调主线连续：博士阶段的半监督学习、多尺度上下文建模和三维分割经验可以迁移到 EM；EM 中的实例级理解、拓扑约束和自动校对方法，也可以反向用于血管、神经、肿瘤和细小解剖结构分析。最后一句可以说：我的研究不是从医学图像简单转向 EM，而是从器官级语义分割拓展到细胞器级实例理解，统一主线是面向少标注复杂三维图像的结构化学习。

\prompt{9. 为什么回国或为什么现在寻找教职？}

希望建立长期、独立的研究方向和学生团队。当前已形成从医学图像到体电镜的连续积累，也具备论文、数据基准、竞赛和工具开发经验，是从博士后合作研究转向独立发展的合适阶段。

\endgroup

\newpage
\sectionbar{Nature 在审合作论文速查}

\prompt{论文信息}

\textbf{题目：}Simultaneous brain-wide single-cell recording resolves spatiotemporal memory architecture

\textbf{公开版本：}\href{https://doi.org/10.64898/2026.05.21.726120}{bioRxiv, 2026.05.21.726120}（2026 年 5 月 22 日发布）；当前为合作论文，投稿 Nature，在审。

\textbf{技术名称：}GLOBE（sinGle-cell spatiotemporaL recOrding Brain-widE）。

\prompt{20 秒说明}

这项工作开发了 GLOBE，通过基因编码的细胞内蛋白磁带记录神经元 Fos 启动子驱动的转录活动，并结合高通量三维图像与信号分析，在单只小鼠全脑范围同时恢复最多 219,703 个神经元连续 5.5 天的活动轨迹，用于研究恐惧学习和记忆提取相关的时空组织规律。

\prompt{一分钟说明}

传统活体记录通常需要在全脑覆盖和单细胞分辨率之间取舍。GLOBE 将随时间变化的转录信号写入神经元内部不断生长的蛋白结构，实验结束后再通过显微成像读取，因此能够在全脑尺度获得单细胞、连续时间轴的转录活动记录。计算端的 GLOBE Tape Reader 完成蛋白磁带和细胞分割、骨架化、四通道信号提取、质量控制、细胞定位及脑图谱配准，再将蛋白磁带上的空间信号转换为时间轨迹。该系统在单只小鼠中记录最多 219,703 个神经元、覆盖 5.5 天，时间戳中位绝对误差为 3.1--6.7 小时。分析发现，与恐惧学习和记忆提取相关的 Fos 转录动态分布于全脑，并呈现脑区特异的时间异质性。

\prompt{必须准确理解的概念}
\begin{itemize}
  \item \textbf{记录对象：}Fos 启动子驱动的转录活动，它是神经元活动的分子代理指标，不是动作电位或毫秒级电活动。
  \item \textbf{读取方式：}实验结束后固定组织并进行显微成像，属于活体记录、离体读取，不是实时全脑显微成像。
  \item \textbf{蛋白磁带：}细胞内自组装且持续生长的蛋白结构；不同荧光单体沿其长度形成空间信号，空间位置可映射为时间。
  \item \textbf{Fos：}神经科学中常用的即时早期基因，可作为神经元受到行为刺激后活动变化的标志。
\end{itemize}

\prompt{核心数字}
\begin{itemize}
  \item 最多同时记录 219,703 个神经元，连续时间跨度为 5.5 天。
  \item 时间戳精度：3.1--6.7 小时中位绝对误差。
  \item 局部记录密度：每个视野中覆盖 69--90\% 的神经元。
  \item 成像读取速度约 2.9 秒/神经元；计算读取速度约 1.0 秒/神经元/GPU。
\end{itemize}

\newpage
\sectionbar{Nature 论文：计算流程与个人贡献}

\begingroup
\small

\prompt{GLOBE Tape Reader 计算流程}
\begin{enumerate}
  \item 输入四通道三维共聚焦图像：Nissl 染色、结构单体、信号单体和时间戳单体。
  \item 使用 MedNeXt-S 进行三维蛋白磁带分割，预测前景、轮廓和骨架距离图，再通过分水岭得到磁带实例。
  \item 将实例转换为三维点云，利用 PCA 引导的样条拟合进行骨架化并生成标准化中心线。
  \item 沿骨架提取四通道强度，通过时间戳信号确定中心位置，计算双侧对称性并完成空间到时间转换。
  \item 使用 micro-SAM 分割神经元胞体，将磁带分配到对应细胞，并进行形态、信号和对称性质量控制。
  \item 将细胞定位并配准到 Allen CCFv3 小鼠脑图谱，完成全脑时空活动分析。
\end{enumerate}

\prompt{计算规模}

单只小鼠最多包含 417 个 ND2 采集文件和 3,038 个图像块，约 7.5 TB 图像数据。分割模型使用 4 张 40 GB NVIDIA A100 训练约 5 小时；单脑完整处理约需 488 GPU 小时和 390 CPU 小时。流程通过 SLURM 作业数组并行执行图像块提取、分割推理和信号分析。

\prompt{论文中的主要发现}
\begin{itemize}
  \item 机器学习识别出 8 类 Fos 转录时间模式，包括编码期激活、持续激活、提取期激活及编码--提取再激活等。
  \item 不同时间模式分布于多个脑区，并呈现明显的脑区特异性。
  \item 全体神经元的时间动态表现出低维结构；随着采样细胞数量增加，对整体结构的恢复更稳定、方差更小。
\end{itemize}

\prompt{你的贡献应如何表述}

论文作者贡献声明明确写明：\textbf{P.L. 参与了蛋白磁带记录计算读取流程的开发，并参与计算分析。}稳妥的口头表述是：

\note{这是一项生物技术、神经科学和计算方法高度交叉的合作工作。我的贡献主要在计算端，参与了蛋白磁带记录的图像读取流程开发和计算分析。整个计算流程需要从大规模三维显微图像中完成蛋白磁带与细胞分割、结构和信号提取、质量控制及并行计算，为后续全脑单细胞时空分析提供数据基础。}

\prompt{交流前必须确认}
\begin{itemize}
  \item 向合作导师或论文主要作者确认自己具体负责的代码模块、实验图或分析结果。
  \item 若不能确认是否直接负责 MedNeXt-S、micro-SAM、信号提取或并行计算中的某一模块，不要将整个 GLOBE Tape Reader 表述为个人开发。
  \item 准确说“投稿 Nature、在审”或“Nature 在审合作论文”，不要说“Nature 论文”或“已被 Nature 接收”。
\end{itemize}

\prompt{可能追问}
\begin{itemize}
  \item \textbf{它与钙成像有什么不同？}GLOBE 可覆盖全脑、单细胞和数天时间跨度，但时间分辨率为小时级且需离体读取；钙成像时间分辨率高，但全脑覆盖与长期单细胞并行记录更困难。
  \item \textbf{主要局限是什么？}Fos 转录只是神经活动的代理指标；时间精度为小时级；当前分析主要是观察性结果，不能直接证明相关细胞群对记忆形成具有因果作用。
  \item \textbf{与你的研究主线有何关系？}计算部分仍是大规模三维生物图像分割、结构提取、质量控制和高通量分析，可自然连接体电镜与医学图像研究。
\end{itemize}

\endgroup

\newpage
\sectionbar{建议向负责人了解的问题}

\begingroup
\small

\prompt{研究与团队}
\begin{itemize}
  \item 团队目前的核心研究方向是什么？最希望新教师补充哪类能力？
  \item 团队通常如何组织合作，青年教师拥有多大的独立发展空间？
  \item 是否已有医院、生命科学团队或其他交叉合作渠道？
  \item 对论文、项目、开源平台和产业合作分别如何评价？
\end{itemize}

\prompt{岗位与考核}
\begin{itemize}
  \item 岗位属于什么类型？聘期和职称晋升路径如何设置？
  \item 前三年的主要考核指标是什么？论文、基金、教学和人才项目分别占多大比重？
  \item 对国家自然科学基金青年项目或人才项目是否有明确预期？
\end{itemize}

\prompt{资源与教学}
\begin{itemize}
  \item 我的研究涉及三维医学图像和体电镜数据，对 GPU 显存、存储及数据读写要求较高。学院或团队目前是否有共享 GPU 服务器或公共算力平台？新教师通常如何申请和使用这些资源？
  \item 团队计算资源以自建服务器还是学校公共平台为主？是否可以使用启动经费购置 GPU 服务器或建设小型计算节点？
  \item 新教师何时能够招收硕士生或博士生？是否有团队共享招生名额？
  \item 第一年的教学工作量如何安排？通常需要承担哪些基础课程？
  \item 是否支持建设数据平台和跨学院联合培养？
\end{itemize}

\prompt{流程}
\begin{itemize}
  \item 本次交流之后通常还有哪些环节和材料要求？
  \item 学院预计的招聘时间表和到岗时间是什么？
\end{itemize}

\note{\textbf{提问策略：}优先问研究匹配和发展空间，再问考核、招生与资源。待遇问题可以了解，但不宜成为第一次交流的主线。若负责人主动谈待遇，应自然、具体地确认岗位类型、合同期限和考核条件。}

\sectionbar{会前与交流中检查表}

\begin{tabularx}{\textwidth}{p{0.7cm} X}
$\square$ & 测试摄像头、麦克风、网络和会议软件，准备备用耳机与手机热点。\\
$\square$ & 打开简历、论文列表和本手册第一页，但不要在回答时持续低头阅读。\\
$\square$ & 准备两篇最熟悉论文的 PDF，以便对方临时追问图表或指标。\\
$\square$ & 查清团队负责人近三年的研究方向和两至三篇代表论文。\\
$\square$ & 准备一个明确的到岗时间范围。\\
$\square$ & 回答先说结论，再说证据；不确定的信息直接说明需要会后确认。\\
$\square$ & 对未发表、在审和拟投稿成果准确区分，不把投稿目标表述为既定结果。\\
$\square$ & 结束时总结匹配点，并确认下一步材料或流程。\\
\end{tabularx}

\endgroup

\sectionbar{个人补充笔记}

\prompt{负责人及团队近期方向：}
\blankline\blankline

\prompt{我与团队最直接的三个结合点：}
\blankline\blankline

\prompt{需要重点展示的成果或资源：}
\blankline\blankline

\prompt{交流后需要跟进的事项：}
\blankline\blankline

\end{document}
